\documentclass[11pt,a4paper]{article}
\usepackage[top=1.6cm,bottom=1.4cm,left=1.3cm,right=1.3cm,headsep=8pt,footskip=16pt]{geometry}
\usepackage{enumitem}
\usepackage{amssymb}
\usepackage{array}
\usepackage[table]{xcolor}
\usepackage{fancyhdr}
\usepackage{multicol}
\usepackage{polyglossia}
\setdefaultlanguage[calendar=gregorian, numerals=maghrib]{arabic}
\setotherlanguage{english}
\IfFontExistsTF{Amiri}{%
  \newfontfamily\arabicfont[Script=Arabic, Scale=1.08]{Amiri}
}{%
  \newfontfamily\arabicfont[Script=Arabic, Scale=1.08, Path=../fonts/, Extension=.ttf,
    UprightFont=*-Regular, BoldFont=*-Bold, ItalicFont=*-Italic,
    BoldItalicFont=*-BoldItalic]{Amiri}
}
\IfFontExistsTF{AlBattar}{%
  \newfontfamily\khatt[Script=Arabic, Scale=1.15]{AlBattar}
}{%
  \newfontfamily\khatt[Script=Arabic, Scale=1.15, Path=../fonts/,
    BoldFont=ae_AlBattar.ttf]{ae_AlBattar.ttf}
}
\renewcommand{\headrulewidth}{0.8pt}
\renewcommand{\footrulewidth}{0.8pt}
\fancyhf{}
\rhead{\khatt المدرسة العليا للأساتذة ببوزريعة}
\lhead{\khatt ملحقة بومرداس}
\lfoot{\khatt الأستاذ: د. بن عبد الرزاق يوسف}
\cfoot{}
\pagestyle{fancy}
\parindent 0pt
\setlength{\columnsep}{0.6cm}
\setlength{\columnseprule}{0.4pt}
\setlist[enumerate,1]{label=\textbf{\arabic*.}, leftmargin=*, itemsep=5pt, topsep=2pt}
\begin{document}
\begin{center}
{\Large\khatt\textbf{مفتاح الإجابات ــ الاختبار في الذكاء الاصطناعي}}\\[4pt]
{\small (نسخة الأستاذ ــ لا تُوزَّع على الطلبة)}\\[10pt]
\renewcommand{\arraystretch}{1.9}\setlength{\tabcolsep}{2.5pt}
\begin{tabular}{|>{\centering\arraybackslash}m{1.9cm}|*{20}{>{\centering\arraybackslash}m{0.5cm}|}}
\hline
\rowcolor{blue!12}\textbf{رقم السؤال} & \textbf{1} & \textbf{2} & \textbf{3} & \textbf{4} & \textbf{5} & \textbf{6} & \textbf{7} & \textbf{8} & \textbf{9} & \textbf{10} & \textbf{11} & \textbf{12} & \textbf{13} & \textbf{14} & \textbf{15} & \textbf{16} & \textbf{17} & \textbf{18} & \textbf{19} & \textbf{20} \\ \hline
\rowcolor{green!16}\textbf{الإجابة} & أ & د & ج & ج & هـ & أ & ج & د & ب & هـ & هـ & ج & د & ب & ب & هـ & ب & أ & د & أ \\ \hline
\end{tabular}

\vspace{10pt}
{\small سلّم التنقيط: نقطة واحدة لكل سؤال (المجموع 20 نقطة).
ترتيب الأسئلة: 1--4 سهلة، 5--10 متوسطة، 11--17 صعبة، 18--20 صعبة جدًا.}
\end{center}
\end{document}
